\documentclass[12pt]{article}
\usepackage[T1]{fontenc}
\usepackage[utf8]{inputenc}
\usepackage{lmodern}
\usepackage{textcomp}
\usepackage{enumitem}
\usepackage{geometry}
\geometry{margin=1in}
\begin{document}
\begin{center}
Answer Key - Astronomy - Graduate Level Assessment 25 \
\small (Answers are ordered exactly as in the question paper.)
\end{center}

\section*{Multiple Choice}
\begin{enumerate}[label=\arabic*.]
\item \textbf{Answer:} B
\\ \textit{Options:} A) The resonance pulls Io in different directions and generates heat.; B) It makes the orbit of Io slightly elliptical.; C) It creates a gap with no asteriods between the orbits.; D) It prevents formation of the ring material into other moons.
\\ \textit{Note:} The 1:2:4 resonance among Jupiter's moons Io, Europa, and Ganymede causes gravitational interactions that lead to Io's orbital eccentricity, making its orbit slightly elliptical.
\item \textbf{Answer:} D
\\ \textit{Options:} A) The Earth must lie completely within the Moon's penumbra.; B) The Moon's penumbra must touch the area where you are located.; C) The Earth must be near aphelion in its orbit of the Sun.; D) The Moon's umbra must touch the area where you are located.
\\ \textit{Note:} To observe a total solar eclipse, it is essential that the Moon's umbra, which is the darkest part of its shadow, reaches the Earth's surface at your location, allowing for the complete obscuration of the Sun.
\item \textbf{Answer:} C
\\ \textit{Options:} A) 150 years; B) 200 years; C) 250 years; D) 300 years
\\ \textit{Note:} The correct answer is 250 years because Pluto has a highly elliptical orbit that takes it significantly longer than the 365 days it takes Earth to complete one orbit around the Sun.
\item \textbf{Answer:} A
\\ \textit{Options:} A) size of the planet; B) presence of an atmosphere; C) distance from the sun; D) rotation period
\\ \textit{Note:} The size of a planet is crucial in determining its volcanic and tectonic activity because larger planets retain heat longer and have sufficient gravitational force to generate tectonic processes, while smaller planets tend to cool quickly and lack the necessary conditions for sustained geological activity.
\item \textbf{Answer:} D
\\ \textit{Options:} A) 2:3 periodic ratio of Neptune:Pluto; B) Kirkwood Gaps.; C) Gaps in Saturn's rings.; D) Breaking of small Jovian moons to form ring materials.
\\ \textit{Note:} The breaking of small Jovian moons to form ring materials is a process primarily driven by tidal forces and collisions rather than orbital resonance, which involves gravitational interactions that maintain stable orbits among larger bodies.
\end{enumerate}

\section*{True / False}
\begin{enumerate}[label=\arabic*.]
\setcounter{enumi}{5}
\item \textbf{Answer:} True
\\ \textit{Options:} A) True; B) False
\\ \textit{Note:} The Cassini division is indeed a significant gap in Saturn's rings that is maintained by the gravitational influence of the moon Mimas, which exerts a tidal force that clears out particles from the region, creating the observable separation.
\item \textbf{Answer:} True
\\ \textit{Options:} A) True; B) False
\\ \textit{Note:} A formation theory of the solar system primarily addresses the processes and conditions that led to the formation of celestial bodies, including planets, but does not require an explanation for the emergence of life, which is a separate scientific inquiry.
\item \textbf{Answer:} True
\\ \textit{Options:} A) True; B) False
\\ \textit{Note:} The statement is true because the Moon is in synchronous rotation, meaning it takes the same amount of time to rotate on its axis as it does to orbit the Earth, which is approximately 27.3 days, leading to the same side of the Moon always facing the Earth.
\item \textbf{Answer:} True
\\ \textit{Options:} A) True; B) False
\\ \textit{Note:} Both Saturn and Jupiter exhibit equatorial wind speeds that do not exceed 900 miles per hour, with Jupiter's maximum equatorial wind speed being around 400 miles per hour and Saturn's being even lower, thus confirming the statement as true.
\item \textbf{Answer:} True
\\ \textit{Options:} A) True; B) False
\\ \textit{Note:} The Arecibo Telescope featured a 305-meter diameter parabolic antenna, making it the largest of its kind in the world until its decommissioning in 2020.
\end{enumerate}

\section*{Long Form Response}
\begin{enumerate}[label=\arabic*.]
\setcounter{enumi}{10}
\item \textbf{Answer:} Referring to Jupiter and other jovian planets as "gas giants" can be misleading, as it overlooks the significant presence of non-gaseous components in their overall composition. While the term emphasizes their vast atmospheres predominantly composed of hydrogen and helium, it fails to account for the substantial amounts of ices, metals, and rocky materials that exist within these planets. For instance, Jupiter's core is believed to consist of a solid mass of heavy elements surrounded by a layer of metallic hydrogen, indicating that its physical structure is more complex than a simple gas giant classification suggests. This nuanced understanding is crucial for comprehending the formation and evolution of these planets, as it highlights the interplay between gaseous and solid materials in their development and the implications for their potential habitability and geophysical processes. Therefore, a more comprehensive terminology is essential for accurately conveying the intricate nature of jovian planets.
\\ \textit{Note:} correct because it highlights that the term "gas giants" inadequately describes jovian planets by neglecting their complex internal structures, which include significant amounts of ices, metals, and rocky materials in addition to their gaseous atmospheres.
\item \textbf{Answer:} Hydrogen plays a pivotal role in the composition and atmospheric dynamics of the Jovian planets, primarily due to its abundance and unique physical properties. As the most prevalent element, hydrogen constitutes the bulk of the atmospheres of planets like Jupiter and Saturn, influencing their overall mass, gravitational forces, and thermal profiles. Its low molecular weight contributes to the formation of deep, convective atmospheres, driving complex weather patterns and the development of phenomena such as storms and jet streams. Compared to heavier elements like helium, ammonia, and methane, hydrogen's dominance dictates the chemical interactions and energy dynamics within these planets, significantly affecting their magnetic fields and interior structures. Thus, hydrogen is not only essential for understanding the atmospheric composition of the Jovian planets but also crucial for elucidating their dynamic behaviors and evolutionary processes.
\\ \textit{Note:} correct because it emphasizes hydrogen's abundance and low molecular weight, which significantly shape the atmospheric composition, dynamics, and weather patterns of the Jovian planets, distinguishing its role from that of other elements.
\end{enumerate}

\vfill
\noindent\textit{End of Answer Key}
\end{document}